\begin{theorem}[Teorema de la Función Implícita]
    Sea \(S \subset \mathbb{R}^n \times \mathbb{R}^m\) abierto y \(f: S \to \mathbb{R}^m\) de clase \(C^1\).
    Sea \((a,b) \in S\) (con \(a \in \mathbb{R}^n, b \in \mathbb{R}^m\)) tal que \(f(a,b) = 0\) y la matriz
    \[ M = \begin{pmatrix}
            \frac{\partial f_1}{\partial x_{n+1}}(a,b) & \dots  & \frac{\partial f_1}{\partial x_{n+m}}(a,b) \\
            \vdots                                     & \ddots & \vdots                                     \\
            \frac{\partial f_m}{\partial x_{n+1}}(a,b) & \dots  & \frac{\partial f_m}{\partial x_{n+m}}(a,b)
        \end{pmatrix} \]
    es inversible.

    Entonces, existen \(A \subset \mathbb{R}^n\) y \(B \subset \mathbb{R}^m\) abiertos tales que \(A \times B \subset S\), \(a \in A\), \(b \in B\), y existe una única función \(g: A \to B\) de clase \(C^1\) tal que \(g(a) = b\) y
    \[ f(x,g(x)) = 0 \quad \forall x \in A \]
\end{theorem}

\noindent\textbf{Observación:} La hipótesis se puede cambiar por pedir que \(Df(a,b)\) sea sobreyectiva, ya que esto obliga a que la matriz tenga \(m\) columnas linealmente independientes. Se puede suponer que son las últimas \(m\), si no se hace un cambio de variables.

\begin{proof}
    Sea \(F: S \subset \mathbb{R}^n \times \mathbb{R}^m \to \mathbb{R}^n \times \mathbb{R}^m\), definida como \(F(x,y) = (x, f(x,y))\). Como \(f \in C^1\), entonces \(F \in C^1\).
    Evaluando en el punto tenemos \(F(a,b) = (a, f(a,b)) = (a,0)\).

    La matriz jacobiana es:
    \[ DF(a,b) = \begin{pmatrix} I & 0 \\ X & M \end{pmatrix} \]
    la cual es inversible porque \(\det(DF(a,b)) = \det(I) \cdot \det(M) = \det(M) \neq 0\).

    Por el Teorema de la Función Inversa aplicado a \(F\), existen \(V \subset S \subset \mathbb{R}^{n+m}\) abierto y \(W \subset \mathbb{R}^{n+m}\) abierto tales que:
    \((a,b) \in V\), \(F(a,b) = (a,0) \in W\), \(F: V \to W\) es biyectiva y \(F^{-1}: W \to V\) es de clase \(C^1\).

    Escribimos \(V = \tilde{A} \times B\) con \(\tilde{A} \subset \mathbb{R}^n\) y \(B \subset \mathbb{R}^m \implies F^{-1}: W \to \tilde{A} \times B\).
    Podemos escribir \(F^{-1}(x,y) = (F_1^{-1}(x,y), F_2^{-1}(x,y))\) con \(F_1^{-1}: W \to \tilde{A}\) y \(F_2^{-1}: W \to B\).

    Aplicando \(F\) se tiene que:
    \[ (x,y) = F(F^{-1}(x,y)) = F(F_1^{-1}(x,y), F_2^{-1}(x,y)) = (F_1^{-1}(x,y), f(F_1^{-1}(x,y), F_2^{-1}(x,y))) \]
    De aquí se deduce que \(F_1^{-1}(x,y) = x \quad \forall (x,y) \in W\) y además se tiene que:
    \[ f(x, F_2^{-1}(x,y)) = y \quad \forall (x,y) \in W \]
    En definitiva, \(F^{-1}(x,y) = (x, F_2^{-1}(x,y))\) con \(F_2^{-1} \in C^1\).

    Notemos que \(\pi_1: \mathbb{R}^n \times \{0\} \to \mathbb{R}^n\), \(\pi_1(x,0) = x\) es de clase \(C^1\) y biyectiva.
    Tomo \(A = \pi_1( (\mathbb{R}^n \times \{0\}) \cap W )\). Como \(W\) es abierto en \(\mathbb{R}^{n+m}\), la intersección es abierta en \(\mathbb{R}^n \times \{0\}\), y por tanto \(A\) es abierto en \(\mathbb{R}^n\).

    Veamos que \(A \subset \tilde{A}\): si \(x \in A \implies x = \pi_1(x,0)\) con \((x,0) \in W\).
    \[ F^{-1}(x,0) = (x, F_2^{-1}(x,0)) \in \tilde{A} \times B \]
    Por lo tanto, \(x \in \tilde{A}\) y \(F_2^{-1}(x,0) \in B\).
    Luego, restringiendo los dominios adecuadamente, \(g: A \to B\) dada por \(g(x) = F_2^{-1}(x,0)\) es de clase \(C^1\).

    Afirmo que \(g(a) = b\). En efecto:
    \[ F(a,b) = (a,0) \implies (a,b) = F^{-1}(a,0) = (a, F_2^{-1}(a,0)) = (a, g(a)) \]
    Por lo tanto, \(g(a) = b\).

    Nos queda ver que \(f(x,g(x)) = 0\).
    \[ f(x, g(x)) = f(x, F_2^{-1}(x,0)) = 0 \]
    pues vimos que \(f(x, F_2^{-1}(x,y)) = y\), y tomando \(y=0\) se cumple.

    Por último, veamos que \(g\) es única.
    Supongamos que \(\exists \tilde{g}: A \to B\) tal que \(f(x, \tilde{g}(x)) = 0 \quad \forall x \in A\).
    \[ F(x, g(x)) = (x, f(x, g(x))) = (x, 0) = (x, f(x, \tilde{g}(x))) = F(x, \tilde{g}(x)) \]
    Como \(F\) es biyectiva en la región, \(F(x, g(x)) = F(x, \tilde{g}(x)) \implies \tilde{g}(x) = g(x) \quad \forall x \in A\).
\end{proof}